\section{Discussion}\label{sec:discussion}

\emph{The score scale is part of the MOT system.} The same detections and trackers give very
different results when only the scale of the scores changes: halving the scores reduced ByteTrack
(official setting) and OC-SORT from 67.70 and 66.43 HOTA to 0, and with RT-DETR-L on UAVDT, ByteTrack
reached a MOTA of only 3.29 because of its many false positives. The fixed thresholds of a tracker
hold only for the detector they were tuned for. The results therefore suggest that the score scale
of the detector is a first-order variable of a tracking-by-detection system and deserves the same
attention as the association method.

\emph{Self-limiting by design.} In the clean mode the layer keeps the tracker's thresholds and
scores, so a tracker with a low-score step and well-matched detections receives its normal input.
This is the main difference from the earlier design \aob{}, which lowered the HOTA of the same
strong trackers by 4.15 to 5.89 points (Table~\ref{tab:prior}). The layer therefore does not need
to know in advance whether a tracker needs help; it infers the operating mode from the score
distribution.

\emph{Tracker dependence.} Single-stage trackers drop every detection below one threshold, so a
track can be lost when its object receives low scores for a few frames. The continuation rescue
fills exactly these gaps, and OC-SORT and the external PD-SORT gained HOTA on 7 of 7 MOT17 sequences.
Two-stage trackers already continue tracks with low-score boxes, so they keep their own accuracy
with well-matched detections and gain when the detector is noisy, as ByteTrack and BoT-SORT did with
RT-DETR-L. The observed gains are thus strongly associated with the host and with how well its
thresholds fit the detector scores.

\emph{Identity switches.} The layer cut identity switches where it removed false boxes: by 43\% and
47\% for ByteTrack and BoT-SORT with RT-DETR-L on KITTI, by 45\% and 56\% on the held-out aerial
data, and by 8 and 14 for OC-SORT on MOT17. Where it let a tracker follow many more objects, the
count rose with the number of followed objects: for OC-SORT with YOLOv8n on KITTI, 79 more switches
came with 5,310 fewer missed objects and 9.60 more IDF1 points. Identities were therefore kept
better overall in this cell too.

\emph{Deployment.} \scl{} needs no training, no GPU, and no access to the inside of the detector or
the tracker. It reads four contract values from the tracker's published configuration and, at run
time, the candidate list, one small grayscale image for the motion cue, and the tracks of the
previous frame. On a CPU it takes 2.95 ms per frame, which is small next to the 42 ms of the
detector, and the tracker itself ran faster with it, from 2.04 to 1.43 ms per frame
(Table~\ref{tab:runtime} and Fig.~\ref{fig:runtime}). The end-to-end time rose by only 4.5\%, and
most of the time per frame is taken by the detector, so the layer adds a low overhead and fits
real-time systems whose detector meets the required frame rate. A new tracker is added through its published
contract values and an adapter, without changes to the rules.
